% Generated draft: requires mathematical and source review.
\documentclass[11pt]{article}
\usepackage[margin=25mm]{geometry}
\usepackage{fontspec}
\setmainfont{Latin Modern Roman}
\usepackage{amsmath,amssymb,mathtools,unicode-math}
\setmathfont{Latin Modern Math}
\usepackage{enumitem,needspace,xurl,hyperref}
\setlength{\parindent}{0pt}
\setlength{\parskip}{5pt}
\setlength{\emergencystretch}{3em}
\newcommand{\N}{\mathbb{N}}
\newcommand{\Z}{\mathbb{Z}}
\newcommand{\Q}{\mathbb{Q}}
\newcommand{\R}{\mathbb{R}}
\newcommand{\C}{\mathbb{C}}
\newcommand{\F}{\mathbb{F}}
\newcommand{\A}{\mathbb{A}}
\newcommand{\PP}{\mathbb{P}}
\newcommand{\E}{\mathbb{E}}
\newcommand{\eps}{\varepsilon}
\newcommand{\dd}{\,\mathrm{d}}
\begin{document}
\section*{UnsolvedMath ID: 20002376}
\textbf{Source problem code:} catalogue entry
\subsection{Definitions and mathematical statement}
James Cummings's contribution consists of four research questions, rather than one fully specified conjecture. For a model $M$ of set theory, $\operatorname{HOD}^M$ is the class, computed in $M$, of hereditarily ordinal-definable sets. The notation $L[\vec E]$ denotes a constructible hierarchy built using an extender sequence, with the fine-structural and coherence requirements to be specified for any particular comparison theorem. The four source questions are:
\begin{enumerate}
\item[(A)] Understand the relationship between models of the form $\operatorname{HOD}^M$ and classical $L[\vec E]$ models.
\item[(B)] Explain exactly how the very strong axioms used in known refutations of UBH and CBH enter those refutations, and whether those axioms are necessary. The original qualification is ``at least in the versions of these results which I have heard about.''
\item[(C)] Explain Woodin's then-recent work on promising extender sequences.
\item[(D)] Describe the status and significance of the $\Omega$ conjecture.
\end{enumerate}
No class of $M$, exact comparison relation in (A), precise versions of the branch hypotheses or strength measure in (B), definition of ``promising'' in (C), or particular formal version of the $\Omega$ conjecture in (D) is given in the supplied extract. In particular, these requests must not be replaced by an invented single implication or an undated assertion about current status.

The research aid addresses only a mechanism relevant to (B). It divides the Neeman--Steel construction into a local-Woodin collision seed, an overlap-lift and ultrapower-identification module, a Martin--Steel well-foundedness input, and a CBH amplification module. For an iteration tree $\mathcal T$ of limit length, write
\[
 \operatorname{WB}(\mathcal T)
 =\{b:b\text{ is a cofinal branch of }\mathcal T
          \text{ whose direct-limit branch model }M_b
          \text{ is well founded}\}.
\]
This definition explains the branch-counting aspect without fixing missing versions of the source hypotheses.

Under the hypotheses of Neeman--Steel Lemma 3.2, the aid reports that equality of the two transported extenders identifies the two branch models with the same ultrapower. Thus, for those branches $b,c$ and transported extenders $E_b,E_c$,
\[
 E_b=E_c\quad\Longrightarrow\quad
 \bigl(M_b\text{ is well founded}\ \Longleftrightarrow\
       M_c\text{ is well founded}\bigr).
\]
Under all the countable-closure hypotheses of their Proposition 2.6, the prescribed lift of a lower tree from $V$ to $\operatorname{Ult}(V,F)$ is reported to preserve its entire well-founded-branch set: the natural branch transport identifies $\operatorname{WB}(\mathcal T)$ with $\operatorname{WB}(\widehat{\mathcal T})$. In particular,
\[
 |\operatorname{WB}(\mathcal T)|\le1
 \quad\Longrightarrow\quad
 |\operatorname{WB}(\widehat{\mathcal T})|\le1
\]
for that construction under those hypotheses. This is an obstruction to producing a uniqueness counterexample by that particular overlap lift. It is not a theorem that all countably closed iteration trees satisfy UBH, and it does not establish the necessity of the original large-cardinal assumptions. The full hypotheses, the precise transport maps and bibliographic identification of Lemma 3.2 and Proposition 2.6 are absent from the supplied record; the displayed claims therefore retain those named conditional qualifications.
\subsection{Short English statement}
Preserve Cummings's four questions about HOD and extender models, branch-hypothesis refutations, promising sequences and the $\Omega$ conjecture. The research aid gives a conditional obstruction for one overlap-lifting mechanism; it does not prove a general branch hypothesis or a lower bound on necessary large-cardinal strength.
\subsection{Sources}
\begin{itemize}
\item James Cummings, contribution A.4, American Institute of Mathematics workshop \emph{Recent advances in core model theory}, \url{https://aimath.org/WWN/coremodel/coremodel.pdf}.
\item Local UnsolvedMath registry, permanent ID 20002376, source code \texttt{catalogue entry}, canonical location \texttt{aim-logic-notes.json}, \texttt{notes[151]}, record dated 14 August 2026. The corrected-verified statement restores the superscript in $\operatorname{HOD}^M$, the vector accent in $L[\vec E]$, and the line-break artifact in ``sequences.''
\item The supplied aid refers to Neeman--Steel Lemma 3.2 and Proposition 2.6, and a Martin--Steel well-foundedness input, without titles, dates or publication details. The original literature field is empty.
\item Supplied research attempt 1: partial result, problem status at run open, rigor labelled proved, novelty confidence low. This modular audit is machine-generated research aid, not peer reviewed; independently verify before citation.
\end{itemize}
\end{document}
